Vision transformers have changed how visual recognition models are designed, but early transformer architectures often scaled poorly with input resolution because self-attention was computed globally over all tokens. The original Swin Transformer addressed this issue by replacing global attention with a hierarchy of local windowed attention blocks and a shifted-window mechanism that lets neighboring windows exchange information without restoring full quadratic attention cost across the entire image \citep{liu2021swin}. This design made transformer-based models much more practical for mainstream computer vision tasks and established Swin as one of the most influential hierarchical transformer backbones.

The limitation that motivated this project is that most implementations and much of the surrounding tooling remain strongly biased toward 2D imagery. In many research and engineering settings, however, the data are not naturally 2D. Medical scans, voxelized CAD objects, and other spatial measurements often live in three dimensions. Video and spatio-temporal signals can require even higher-dimensional reasoning. Reusing the shifted-window idea in these settings should be possible in principle, but real codebases often encode 2D assumptions in patch extraction, window bookkeeping, relative position indexing, data loading, and training orchestration. As a result, extending a 2D Swin implementation into a general N-dimensional system is less a matter of changing a single attention layer than of making the entire modeling and experimentation stack dimension-aware.

The goal of this practical work was therefore to build an N-dimensional shifted-window transformer in JAX and to verify that the resulting implementation can be trained in a realistic, traceable workflow. The emphasis of the project is not on proposing a new transformer variant. Instead, the contribution is an implementation and systems contribution: one codebase, one configuration-driven command surface, and one training stack that can support both 2D image classification and 3D voxel classification by changing configuration rather than rewriting architecture-specific code. This project also aimed to make the surrounding workflow credible from a research-practice perspective by supporting hyperparameter sweeps, queue-based job execution, logging, checkpointing, and dataset validation before a run starts.

The repository snapshot analyzed in this report contains exactly the kind of evidence needed for such a claim. It includes preserved sweeps and restored-best retraining checkpoints for CIFAR-10, CIFAR-100, and a voxelized ModelNet40 dataset. The final checkpoints reach validation top-1 accuracies of \(0.9182\), \(0.6948\), and \(0.7421\), respectively, with recomputed held-out test top-1 accuracies of \(0.8995\), \(0.6963\), and \(0.6965\). It also includes queue, tmux, and resume logs showing how the workloads were orchestrated through the same CLI-based pipeline, including the ModelNet40 rerun after a train-derived validation split was created. These artifacts demonstrate that the project is not merely a collection of isolated modules, but a functioning experimentation framework.

This report is structured as follows. \autoref{sec:related-work} situates the implementation in the context of the original Swin paper, the JAX/Flax ecosystem, and the datasets used here. \autoref{sec:methods} explains how the architecture and training system were generalized from fixed 2D assumptions to arbitrary spatial dimensionalities. \autoref{sec:setup} documents the experimental setup, the datasets, the hardware, and the sweep procedures. \autoref{sec:results} interprets the validation-selected checkpoints and the recomputed test metrics, including the gap between sweep-best and final retraining performance for ModelNet40. Finally, \autoref{sec:limitations} summarizes the main constraints of the preserved evidence, especially the lack of repeated-seed uncertainty estimates. The report deliberately stays within those limits so that every conclusion can be traced back to the saved artifacts rather than to assumptions or reconstructed claims.
